\section{Introduction}
\label{sec:intro}

On 28 February 2026 the United States and Israel began sustained strikes on Iran, and within weeks traffic through the Strait of Hormuz was interrupted at intervals. Brent peaked some \$42 a barrel, or 57 per cent, above its pre-war level---comparable to the \$35 increase of 2022---while UK natural gas peaked 78 pence per therm above pre-war, roughly a quarter of the 2022 move. Petrol rose 20 per cent and diesel 36 per cent. Both the IMF and the OECD cut their 2026 UK growth forecasts by half a percentage point, the largest markdown applied to any G7 economy.

That the UK should take the sharpest hit follows from three structural features. Gas is 62 per cent of final household energy consumption \citep{desnz2025ecuk}, the highest share among the European members of the G7 --- Italy, next highest, is at 53 per cent; no published source ranks household gas shares across the G7 as a whole, the comparable Eurostat series excluding the United States, Canada and Japan \citep{eurostat2018}, so the claim is stated at that strength. Gas-fired plant is the marginal generator roughly 85 per cent of the time, so the same wholesale move arrives a second time through the bills of households that use no gas at all. And the retail price cap transmits wholesale prices mechanically but with a lag of roughly one and a half quarters, so a spring 2026 shock lands on the winter 2026--27 heating season and on the fiscal year the Autumn Budget will set: Ofgem confirmed an October cap of \pounds 1{,}723 on 26 August 2026 \citep{ofgem2026cap}, on the reduced typical-consumption basis effective 1 July 2026 \citep{ofgem2026tdcv}.

The policy question is live. The Chancellor has ruled out universal support of the Energy Price Guarantee kind and, as trailed in press reporting rather than in any published Treasury document, is said to favour an income-based social tariff; I treat that as a reported signal rather than as policy. The Joseph Rowntree Foundation, joined by the TUC, counters with a universal discounted block \citep{jrf2026} costing near \pounds 5 billion. The statistic anchoring every Budget submission is that around 40 per cent of households struggling to heat their homes are not on means-tested benefits: if it is right, the distinction between the proposals is coverage rather than generosity. The Warm Home Discount reform reported for August 2026 --- \pounds 150, automatic in England and Wales, said to add 2.7 million households --- would move towards coverage without settling the question; here too no published Government document sets it out in those terms, and it is modelled as reported.

Answering the question requires evidence that does not exist for the UK. The macroeconomic forecasters produce scenario paths for oil, gas and inflation but stop at the aggregate; the distributional analysts produce household estimates but take the price path as a stipulated assumption. This paper maps a macro oil-and-gas scenario through the Ofgem cap formula and pump prices into a household price vector, runs that vector through UK microdata to give first-order incidence by income decile, and scores five instruments from the Autumn Budget debate against the resulting loss.

\paragraph{Identification: the pre-war cap is solved, not observed.} Everything turns on what a household would have paid in the twelve months from March 2026 had the war not happened, and that cannot be read off a published cap: Ofgem's cap prices a forward window closing about seven weeks before the charge period, so the window behind the July 2026 cap of \pounds 1{,}663 already contained the war. Treating it as the pre-shock baseline while assigning the same quarter full war pass-through subtracts part of the shock from itself, understating the domestic leg at source and inflating the motor-fuel share by exactly what the domestic leg loses. The pre-war level is therefore solved jointly with the sustained fraction of the wholesale gas peak from both published caps --- two equations, two unknowns --- giving a pre-war counterfactual cap of \pounds \genCapBaselineLevel{}. That it reproduces both caps to \pounds \genCapFitMaxErrorGbp{} is not evidence for it, zero degrees of freedom guaranteeing a zero residual; what the solve should be judged on is its fragility to the assumed monthly shape of the gas peak, the one input the formula consumes that is not observed at the frequency it needs. Sweeping that profile moves the solved cap from \pounds \genGasProfileCapMinGbp{} to \pounds \genGasProfileCapMaxGbp{} and the aggregate loss from \pounds \genGasProfileAggMinBn{} to \pounds \genGasProfileAggMaxBn{} billion, with \genGasProfileNonIdentifiedCount{} of \genGasProfileCellCount{} cells admitting no solution, and the central case sits near the bottom of the swept cap range and near the top of the swept aggregate range (Section~\ref{sec:step1}, Appendix~\ref{app:domesticleg}).

\paragraph{Contributions.} Three specification choices govern everything downstream. Both legs of the shock are annualised over one \emph{window}, March 2026 to February 2027, with calendar 2026 retained separately so the Resolution Foundation's \pounds 11 billion \citep{rf2026outlook} can be compared on its own window. The \emph{accounting basis} is the consumption-weighted average of the quarterly cap levels over that window, the steady-state level being a labelled variant. The \emph{denominator} is equivalised household net income after housing costs on the HBAI scale.

The mean annual household loss is \pounds \genCentralLossMean{}, or \genCentralLossPctIncome{} per cent of equivalised income, aggregating to \pounds \genAggSpendRealisedBn{} billion; on the steady-state basis, \pounds \genSteadyStateLossMean{} and \pounds \genSteadyStateAggBn{} billion. Re-run on calendar 2026, the damped and peak-fuel specifications give \pounds \genCalendarAggBn{} to \pounds \genPeakFuelCalendarAggBn{} billion, bracketing the Resolution Foundation's estimate on a single window.

Motor fuel carries \genMotorFuelShareOfLoss{} per cent of the aggregate loss, against \genGasShareOfLoss{} per cent through gas and \genElecShareOfLoss{} per cent through electricity: the largest single channel, but a bare majority rather than a dominant one, and \genSpecFuelShareMinPct{} to \genSpecFuelShareMaxPct{} per cent across the seven specifications of Section~\ref{sec:results-fuel}. The claim defended is the one that holds across the range: a substantial part of the loss, never below about a third, runs through a channel that every instrument in the debate --- social tariff, discounted block, Warm Home Discount, VAT zero-rating, all of them domestic-bill instruments --- leaves untouched.

The distributional contribution crosses two facts the literature establishes separately. \citet{fetzer2024} show that affluent English neighbourhoods are more exposed in \emph{pounds}; the budget-share literature shows that poorer households lose more in \emph{per cent}. Displayed on the same households, the percentage gradient is steep --- \genDecileOneLossPct{} per cent in decile one against \genDecileTenLossPct{} per cent in decile ten --- while the cash loss is not monotone: the largest falls on decile eight, at \pounds \genDecileEightLossGbp{}, above decile ten's \pounds \genDecileTenLossGbp{}. Vehicle ownership and mileage, not dwelling efficiency, drive the cash profile. The direction of the percentage gradient is robust to every specification here; its magnitude varies by more than a factor of two, and no claim rests on which measurement a reader prefers.

Finally, an exchequer envelope buys more on eligibility than on generosity. Held at their feasible maxima the social tariff and the Warm Home Discount both exhaust at \pounds \genSocialTariffFeasibleMaxCostBn{} billion --- the same figure by construction --- and offset \genSocialTariffEnvelopeOffsetPct{} and \genWhdEnvelopeOffsetPct{} per cent of the aggregate loss, while spending on eligibility at each sponsor's own generosity offsets \genSocialTariffEligibilityOffsetPct{} and \genWhdEligibilityOffsetPct{} per cent, at \pounds \genSocialTariffEligibilityCostBn{} and \pounds \genWhdEligibilityCostBn{} billion. That is not the same money --- the eligibility arm spends roughly a quarter more --- so the comparison is stated per pound in Section~\ref{sec:policy-eligibility}, and it is conditional on an imputation this paper cannot validate: under a parity fuel specification the advantage falls to \genTargetingMultipleUnderParity{} times, and the range rather than the point is what should be quoted.

\paragraph{What this paper does not do.} It contains no general-equilibrium analysis. \citet{kanzig2023} finds that general-equilibrium and indirect effects account for roughly two-thirds of the household response to an energy shock, so a static consumption-side microsimulation captures at most a third of true incidence; \citet{goulder2019} add that it observes only the use side, and \citet{pallotti2024} that euro-area heterogeneity in 2021--23 was driven by nominal net asset positions rather than the energy basket. I do not model consumption substitution and do not deduct baseline VAT from net income. Each is stated at the point of use and revisited in Section~\ref{sec:discussion}.

Section~\ref{sec:literature} locates the paper in the literature. Section~\ref{sec:methodology} sets out the price mapping and the microsimulation. Section~\ref{sec:results} reports the price path, incidence and the fuel decomposition. Section~\ref{sec:policy} scores the five proposals. Section~\ref{sec:discussion} states the limitations and the agenda they imply.
